
\chapter{卷积神经网络——让机器看懂世界}

\intro{卷积神经网络（CNN）是计算机视觉领域的基石技术，它在图像分类、目标检测、图像分割等任务上的性能已超越人类水平。本章从卷积的数学定义出发，深入理解CNN的每一个核心组件——卷积、池化、通道交互——并分别用NumPy从零实现卷积和PyTorch构建完整的CNN训练管线。本章包含大量Python代码示例和可视化，是全书代码最密集的一章。}

%%%%%%%%%%%%%%%%%%%%%%%%%%%%%%%%%%%%%%%%%%%%%%%%%%%%%%%%
\section{视觉问题的本质}

\subsection{全连接网络的失效}

考虑一张 $224 \times 224 \times 3$ 的RGB图像（ImageNet标准尺寸）。如果将其展平为150528维的向量并输入一个全连接网络，第一隐藏层假设有1000个神经元：

\[
\text{第一层权重参数量} = 150528 \times 1000 \approx 1.5 \times 10^8
\]

仅第一层就有1.5亿个参数。这导致三个致命问题：
\begin{enumerate}
    \item \textbf{计算不可行}——每次前向传播需要1.5亿次乘法
    \item \textbf{严重过拟合}——参数量远超样本量
    \item \textbf{忽略空间结构}——将图像展平为向量丢失了像素之间的空间关系
\end{enumerate}

\subsection{图像的三个归纳偏置}

CNN的设计基于对图像数据的三个核心归纳偏置（inductive biases）：

\begin{itemize}
    \item \textbf{局部性}（Locality）——图像的语义信息存在于邻近像素构成的局部区域中，远处的像素通常与当前像素无关
    \item \textbf{平移等变性}（Translation Equivariance）——无论一只猫出现在图像的左上角还是右下角，用于检测猫的滤波器应该在整张图上都有效
    \item \textbf{层次化组合}（Hierarchical Composition）——边缘组成纹理，纹理组成部件，部件组成物体
\end{itemize}

\begin{figure}[H]\centering
\begin{tikzpicture}[scale=0.85]
\node[font=\small\bfseries] at (-6, 2.0) {全连接};
\node[font=\small\bfseries] at (-2, 2.0) {卷积（局部）};

% FC version
\foreach \i in {1,...,6} {
    \foreach \j in {1,...,6} {
        \fill[blue!10] (\i*0.4 -6.5, -0.2-\j*0.4) rectangle ++(0.35,0.35);
    }
}
\draw[->, very thick, red!60] (-5.3, -1.0) -- (-0.5, -0.5);
\node[font=\footnotesize, red!70] at (-2.5,-2.0) {全连接：每像素 \quad 都与下一层所有神经元连接};

% Conv version
\foreach \i in {1,...,6} {
    \foreach \j in {1,...,6} {
        \fill[blue!10] (\i*0.4 -2.5, -0.2-\j*0.4) rectangle ++(0.35,0.35);
    }
}
\draw[red!60, very thick] (0.3, -0.1) rectangle (1.0, -1.5);
\draw[->, very thick, red!60] (0.65, -1.5) -- (2.0, -2.5);
\fill[red!20] (2.0, -3.0) rectangle (2.35, -3.35);
\node[font=\footnotesize, red!70] at (2.2,-3.8) {卷积：仅局部邻域 \quad 连接到下一层};

\node[font=\footnotesize] at (-1.0,-4.2) {参数从 $O(N^2)$ 降至 $O(K^2)$，$N$为图像尺寸，$K$为卷积核大小};
\end{tikzpicture}
\caption{全连接（左）将每个像素独立连接到所有隐藏神经元，参数爆炸且丢失空间结构。卷积（右）仅处理局部邻域，参数共享，保留空间信息}
\label{fig:fc-vs-conv}
\end{figure}

%%%%%%%%%%%%%%%%%%%%%%%%%%%%%%%%%%%%%%%%%%%%%%%%%%%%%%%%
\section{卷积运算的数学定义}

\subsection{一维卷积}

给定两个函数 $f$ 和 $g$，它们的卷积定义为：

\[
(f * g)(t) = \int_{-\infty}^{\infty} f(\tau)\,g(t - \tau)\,d\tau
\]

在离散域中（深度学习均为离散情况）：

\[
(f * g)[n] = \sum_{m=-\infty}^{\infty} f[m]\,g[n - m]
\]

在深度学习中，我们使用的是\textbf{互相关}（cross-correlation），但习惯上仍称为``卷积''。两者的区别仅在于卷积核是否翻转180度——由于卷积核是可学习的参数，翻转与否不影响最终学习结果。

\begin{mydefinition}{卷积与互相关}
深度学习中实际使用的``卷积''操作是\textbf{互相关}：$(f \star k)[i] = \sum_{m} f[i+m]\,k[m]$。它与数学卷积的区别是 $k$ 不翻转。因为 $k$ 由梯度下降学习得到，翻转与否无关紧要。
\end{mydefinition}

\subsection{二维卷积——图像处理的核心}

对于输入图像 $\mathbf{X} \in \R^{H \times W}$ 和卷积核 $\mathbf{K} \in \R^{k_h \times k_w}$，二维卷积（互相关）的输出为：

\[
\mathbf{Y}[i, j] = \sum_{m=0}^{k_h-1}\sum_{n=0}^{k_w-1} \mathbf{X}[i+m,\, j+n] \cdot \mathbf{K}[m, n]
\]

\begin{figure}[H]\centering
\begin{tikzpicture}[scale=1.0]
% Input matrix
\foreach \i in {0,...,4} {
    \foreach \j in {0,...,4} {
        \pgfmathsetmacro{\val}{int((\i+\j)*2+1)}
        \fill[blue!5, draw=gray] (\i*0.8, 3.0-\j*0.8) rectangle ++(0.8,0.8);
        \node[font=\tiny] at (\i*0.8+0.4, 3.0-\j*0.8+0.4) {$\val$};
    }
}
\node[font=\footnotesize, below=0.1cm] at (1.6,-0.8) {输入 $\mathbf{X}$，$5\times5$};

% Kernel
\foreach \i in {0,...,2} {
    \foreach \j in {0,...,2} {
        \pgfmathsetmacro{\kval}{int(\i+\j+1)}
        \fill[red!15, draw=red!60] (\i*0.8+5.5, 3.0-\j*0.8) rectangle ++(0.8,0.8);
        \node[font=\tiny] at (\i*0.8+5.9, 3.0-\j*0.8+0.4) {$\kval$};
    }
}
\node[font=\footnotesize] at (6.7,-0.8) {卷积核 $\mathbf{K}$，$3\times3$};

% Arrow
\draw[->, very thick] (4.3,1.2) -- (5.5,1.2) node[midway, above, font=\footnotesize] {$*$};

% Output
\foreach \i in {0,...,2} {
    \foreach \j in {0,...,2} {
        \fill[green!15, draw=green!60!black] (\i*0.8+9.5, 3.0-\j*0.8) rectangle ++(0.8,0.8);
        \node[font=\tiny] at (\i*0.8+9.9, 3.0-\j*0.8+0.4) {$?$};
    }
}
\node[font=\footnotesize] at (10.7,-0.8) {输出 $\mathbf{Y}$，$3\times3$};

% Arrows illustrating sliding
\draw[blue!60, dashed, ->] (0.4,3.8) -- (2.8,3.8) node[midway, above, font=\tiny] {滑动窗口};
\draw[blue!60, dashed, ->] (4.4,3.5) -- (4.4,1.1) node[midway, left, font=\tiny] {逐元素};
\draw[blue!60, dashed, ->] (4.4,0.7) -- (4.4,-0.1) node[midway, left, font=\tiny] {求和};

\node[font=\footnotesize] at (-0.2,-1.2) {$i$};
\node[font=\footnotesize, rotate=90] at (-0.5,3.0) {$j$};
\end{tikzpicture}
\caption{二维卷积的可视化。$3\times3$卷积核在$5\times5$输入上以步长1滑动，每个位置计算核与局部区域的内积，得到$3\times3$的特征图}
\label{fig:conv2d-illustration}
\end{figure}

\subsection{多通道卷积}

真实图像具有多个通道（如RGB三通道）。对于输入 $\mathbf{X} \in \R^{C_{\text{in}} \times H \times W}$ 和卷积核 $\mathbf{K} \in \R^{C_{\text{out}} \times C_{\text{in}} \times k_h \times k_w}$，多通道卷积为：

\[
\mathbf{Y}[o, i, j] = \sum_{c=0}^{C_{\text{in}}-1}\sum_{m=0}^{k_h-1}\sum_{n=0}^{k_w-1} \mathbf{X}[c,\, i+m,\, j+n] \cdot \mathbf{K}[o, c, m, n] + \mathbf{b}[o]
\]

其中 $o$ 是输出通道的索引。直观理解：每个输出通道的卷积核是一组（$C_{\text{in}}$个）$k_h \times k_w$ 的滤波器，分别作用于输入的每个通道后求和，再加上该输出通道的偏置。

\subsection{步长与填充}

\textbf{步长}（stride）控制卷积核在输入上的滑动步距：
\[
H_{\text{out}} = \left\lfloor\frac{H_{\text{in}} - k_h}{s} + 1\right\rfloor
\]

\textbf{填充}（padding）在输入的四周添加额外的行/列（通常填0），以控制输出尺寸：

\[
H_{\text{out}} = \left\lfloor\frac{H_{\text{in}} + 2p - k_h}{s} + 1\right\rfloor
\]

当 $s = 1$ 且 $p = (k-1)/2$ 时（kernel\_size为奇数），输出尺寸与输入相同——称为 \textbf{``same'' padding}。这是现代CNN设计中的标准配置（如ResNet中的 $3\times3$ 卷积）。

\begin{figure}[H]\centering
\begin{tikzpicture}[scale=0.85]
% Without padding
\node[font=\small\bfseries] at (-4, 3.0) {无填充，步长2};

\foreach \i in {0,...,5} {
    \foreach \j in {0,...,5} {
        \fill[blue!5, draw=gray, thin] (\i*0.6 -4.5, 2.5-\j*0.6) rectangle ++(0.6,0.6);
    }
}
\draw[red!70, very thick] (-4.5,2.5) rectangle (-2.7,0.7);
\draw[->, red!70, thick] (-2.5,1.6) -- (-1.0,1.6);
\foreach \i in {0,...,1} {
    \foreach \j in {0,...,1} {
        \fill[green!15, draw=green!60!black, thin] (\i*0.6 -0.5, 2.5-\j*0.6) rectangle ++(0.6,0.6);
    }
}
\node[font=\footnotesize] at (-2.1,-0.8) {$6\times6 \to 2\times2$};

% With padding
\node[font=\small\bfseries] at (4.0, 3.0) {Same Padding, 步长1};

\foreach \i in {0,...,5} {
    \foreach \j in {0,...,5} {
        \fill[blue!5, draw=gray, thin] (\i*0.6+3.5, 2.5-\j*0.6) rectangle ++(0.6,0.6);
    }
}
\draw[red!70, very thick] (2.9,1.9) rectangle (4.7,0.1);
\draw[->, red!70, thick] (4.9,1.0) -- (6.4,1.0);
\foreach \i in {0,...,5} {
    \foreach \j in {0,...,5} {
        \fill[green!15, draw=green!60!black, thin] (\i*0.6+7.0, 2.5-\j*0.6) rectangle ++(0.6,0.6);
    }
}
\node[font=\footnotesize] at (8.5,-0.8) {$6\times6 \to 6\times6$};

\node[font=\footnotesize] at (1.5,-1.5) {输出尺寸 = $\lfloor\frac{\text{输入} + 2\times\text{填充} - \text{核大小}}{\text{步长}}\rfloor + 1$};
\end{tikzpicture}
\caption{步长和填充对输出尺寸的影响。Same Padding使输入输出尺寸保持一致}
\label{fig:stride-padding}
\end{figure}

%%%%%%%%%%%%%%%%%%%%%%%%%%%%%%%%%%%%%%%%%%%%%%%%%%%%%%%%
\section{用 NumPy 从零实现二维卷积}

理解卷积的最直接方式是亲手实现它。以下代码不使用任何深度学习框架，纯NumPy实现多通道二维卷积：

\begin{mycode}{NumPy 实现完整的 conv2d}
\begin{lstlisting}[language=Python]
import numpy as np

def conv2d_numpy(x, kernel, bias=None, stride=1, padding=0):
    """
    纯NumPy二维卷积（互相关）实现。

    参数:
        x: 输入，shape (C_in, H, W)
        kernel: 卷积核，shape (C_out, C_in, kH, kW)
        bias: 偏置，shape (C_out,) 或 None
        stride: 步长（int）
        padding: 填充（int）

    返回:
        output: shape (C_out, H_out, W_out)
    """
    C_in, H, W = x.shape
    C_out, C_in_k, kH, kW = kernel.shape
    assert C_in == C_in_k, "Input channels must match kernel channels"

    # 添加零填充
    if padding > 0:
        x_padded = np.pad(x, ((0,0), (padding,padding),
                              (padding,padding)),
                          mode='constant', constant_values=0)
    else:
        x_padded = x

    # 计算输出尺寸
    H_out = (H + 2 * padding - kH) // stride + 1
    W_out = (W + 2 * padding - kW) // stride + 1
    output = np.zeros((C_out, H_out, W_out))

    # 卷积循环（理解原理用，生产环境用im2col或FFT）
    for o in range(C_out):
        for i in range(H_out):
            for j in range(W_out):
                h_start = i * stride
                w_start = j * stride
                # 提取局部区�?                patch = x_padded[:,
                         h_start:h_start + kH,
                         w_start:w_start + kW]  # (C_in, kH, kW)
                # 内积 + 偏置
                output[o, i, j] = np.sum(
                    patch * kernel[o]
                ) + (bias[o] if bias is not None else 0)

    return output

# ========== 测试：边缘检测 ==========
# 创建测试图像（6x6渐变）
x = np.arange(36, dtype=np.float32).reshape(1, 6, 6)
print("输入图像 (1, 6, 6):\n", x[0])

# Sobel 水平边缘检测核
sobel_x = np.array([[[[-1, 0, 1],
                      [-2, 0, 2],
                      [-1, 0, 1]]]], dtype=np.float32)
# 输出通道数 = 1

out = conv2d_numpy(x, sobel_x, stride=1, padding=0)
print("\n水平边缘检测结果 (1, 4, 4):\n", out[0])

# ========== 验证：与PyTorch对比 ==========
import torch
import torch.nn.functional as F

x_t = torch.tensor(x).unsqueeze(0)     # (1, C_in, H, W)
k_t = torch.tensor(sobel_x)             # (C_out, C_in, kH, kW)
b_t = torch.zeros(k_t.shape[0])

out_torch = F.conv2d(x_t, k_t, bias=b_t, stride=1, padding=0)
print("\nPyTorch 输出:\n", out_torch.squeeze().numpy())

diff = np.abs(out - out_torch.squeeze().numpy()).max()
print(f"\n最大误差: {diff:.2e}  (应当 < 1e-6)")
\end{lstlisting}
\end{mycode}

\subsection{卷积的矩阵乘法实现——im2col}

上述三层嵌套循环便于理解但效率极低。现代卷积实现使用 \textbf{im2col} 技巧将卷积转化为高效的矩阵乘法：

\begin{mycode}{im2col —— 将卷积转化为矩阵乘法}
\begin{lstlisting}[language=Python]
import numpy as np

def im2col(x, kH, kW, stride=1, padding=0):
    """
    将图像展开为列矩阵，每个局部窗口变为一个列向量。

    参数:
        x: (C_in, H, W)
    返回:
        cols: (C_in * kH * kW, H_out * W_out)
    """
    C_in, H, W = x.shape

    if padding > 0:
        x = np.pad(x, ((0,0), (padding,padding),
                       (padding,padding)),
                   mode='constant')

    H_out = (H + 2*padding - kH) // stride + 1
    W_out = (W + 2*padding - kW) // stride + 1

    cols = np.zeros((C_in * kH * kW, H_out * W_out))

    col_idx = 0
    for i in range(0, H_out * stride, stride):
        for j in range(0, W_out * stride, stride):
            patch = x[:, i:i+kH, j:j+kW]           # (C_in, kH, kW)
            cols[:, col_idx] = patch.reshape(-1)
            col_idx += 1

    return cols, H_out, W_out

def conv2d_im2col(x, kernel, bias=None, stride=1, padding=0):
    """
    使用 im2col 的快速卷�?    """
    C_in, H, W = x.shape
    C_out, C_in_k, kH, kW = kernel.shape

    # 将卷积核展平为 (C_out, C_in*kH*kW)
    kernel_flat = kernel.reshape(C_out, -1)

    # im2col
    cols, H_out, W_out = im2col(x, kH, kW, stride, padding)

    # 矩阵乘法: (C_out, C_in*kH*kW) @ (C_in*kH*kW, H_out*W_out)
    out_flat = kernel_flat @ cols               # (C_out, H_out*W_out)

    output = out_flat.reshape(C_out, H_out, W_out)
    if bias is not None:
        output += bias.reshape(-1, 1, 1)
    return output

# ========== 性能对比 ==========
import time

x_large = np.random.randn(3, 224, 224)
kernel_large = np.random.randn(64, 3, 3, 3)

# Naive 实现
t0 = time.time()
out_naive = conv2d_numpy(x_large, kernel_large, stride=1)
t_naive = time.time() - t0
print(f"Naive: {t_naive:.2f}s")

# im2col 实现
t0 = time.time()
out_im2col = conv2d_im2col(x_large, kernel_large, stride=1)
t_im2col = time.time() - t0
print(f"im2col: {t_im2col:.2f}s")
print(f"加速比: {t_naive/t_im2col:.1f}x")
print(f"结果一致: {np.allclose(out_naive, out_im2col, atol=1e-5)}")
\end{lstlisting}
\end{mycode}

%%%%%%%%%%%%%%%%%%%%%%%%%%%%%%%%%%%%%%%%%%%%%%%%%%%%%%%%
\section{池化层——降采样与平移不变性}

池化（Pooling）层在特征图的每个局部窗口内执行聚合操作，降低空间分辨率的同时保留重要信息。

\subsection{最大池化与平均池化}

\[
\text{MaxPool: } \mathbf{Y}[i, j] = \max_{m,n \in \text{window}} \mathbf{X}[i \cdot s + m,\; j \cdot s + n]
\]

\[
\text{AvgPool: } \mathbf{Y}[i, j] = \frac{1}{k^2}\sum_{m,n \in \text{window}} \mathbf{X}[i \cdot s + m,\; j \cdot s + n]
\]

池化的作用：
\begin{itemize}
    \item \textbf{降采样}——将特征图的空间尺寸缩小，减少后续层的计算量
    \item \textbf{增大感受野}——同样的卷积核在更小的特征图上覆盖相对更大的输入范围
    \item \textbf{局部平移不变性}——小的平移不会改变池化输出（MaxPool对小扰动鲁棒）
    \item \textbf{特征压缩}——保留最显著的特征响应，丢弃冗余信息
\end{itemize}

\begin{mycode}{NumPy 实现最大池化}
\begin{lstlisting}[language=Python]
import numpy as np

def maxpool2d_numpy(x, kernel_size=2, stride=None):
    """
    纯NumPy最大池化

    参数:
        x: (C, H, W)
        kernel_size: 池化窗口大小
        stride: 步长（默认 = kernel_size）
    """
    if stride is None:
        stride = kernel_size

    C, H, W = x.shape
    H_out = (H - kernel_size) // stride + 1
    W_out = (W - kernel_size) // stride + 1

    output = np.zeros((C, H_out, W_out))

    for c in range(C):
        for i in range(H_out):
            for j in range(W_out):
                h_start = i * stride
                w_start = j * stride
                window = x[c,
                          h_start:h_start + kernel_size,
                          w_start:w_start + kernel_size]
                output[c, i, j] = window.max()

    return output

# 测试
feat_map = np.array([[
    [1, 3, 2, 4],
    [5, 6, 8, 7],
    [3, 2, 1, 0],
    [9, 5, 4, 2]
]], dtype=np.float32)

print("输入特征�? (1,4,4):\n", feat_map[0])
print("\nMaxPool (2x2, stride=2):\n",
      maxpool2d_numpy(feat_map, 2, 2)[0])

# PyTorch 验证
import torch
import torch.nn.functional as F
x_t = torch.tensor(feat_map).unsqueeze(0)
out_torch = F.max_pool2d(x_t, kernel_size=2, stride=2)
print("\nPyTorch 输出:\n", out_torch.squeeze().numpy())
\end{lstlisting}
\end{mycode}

%%%%%%%%%%%%%%%%%%%%%%%%%%%%%%%%%%%%%%%%%%%%%%%%%%%%%%%%
\section{PyTorch CNN 完整训练管线}

下面的代码展示了一个完整的CNN分类器——从数据加载到训练、评估和推理——用于CIFAR-10图像分类。

\begin{mycode}{完整的 PyTorch CNN 训练管线（CIFAR-10）}
\begin{lstlisting}[language=Python]
import torch
import torch.nn as nn
import torch.nn.functional as F
import torch.optim as optim
from torch.utils.data import DataLoader
import torchvision
import torchvision.transforms as transforms
import matplotlib.pyplot as plt
import numpy as np
from tqdm import tqdm

# ========== 数据增强与标准化 ==========
train_transform = transforms.Compose([
    transforms.RandomCrop(32, padding=4),     # 随机裁�?    transforms.RandomHorizontalFlip(),         # 随机水平翻转
    transforms.ColorJitter(0.1, 0.1, 0.1),    # 颜色抖动
    transforms.ToTensor(),
    transforms.Normalize(
        (0.4914, 0.4822, 0.4465),             # CIFAR-10 均值
        (0.2470, 0.2435, 0.2616)              # CIFAR-10 标准�?    )
])

test_transform = transforms.Compose([
    transforms.ToTensor(),
    transforms.Normalize(
        (0.4914, 0.4822, 0.4465),
        (0.2470, 0.2435, 0.2616)
    )
])

train_data = torchvision.datasets.CIFAR10(
    root='./data', train=True,
    download=True, transform=train_transform
)
test_data = torchvision.datasets.CIFAR10(
    root='./data', train=False,
    download=True, transform=test_transform
)

train_loader = DataLoader(train_data, batch_size=128,
                          shuffle=True, num_workers=2,
                          pin_memory=True)
test_loader = DataLoader(test_data, batch_size=100,
                         shuffle=False, num_workers=2)

# CIFAR-10 类别
classes = ('plane', 'car', 'bird', 'cat', 'deer',
           'dog', 'frog', 'horse', 'ship', 'truck')

# ========== CNN 模型定义 ==========
class CIFAR10CNN(nn.Module):
    def __init__(self, num_classes=10):
        super().__init__()
        # Block 1: 3 -> 32 通道
        self.conv1 = nn.Conv2d(3, 32, kernel_size=3,
                               padding=1, bias=False)
        self.bn1 = nn.BatchNorm2d(32)
        self.conv2 = nn.Conv2d(32, 32, kernel_size=3,
                               padding=1, bias=False)
        self.bn2 = nn.BatchNorm2d(32)
        self.pool1 = nn.MaxPool2d(2, 2)
        self.dropout1 = nn.Dropout2d(0.1)

        # Block 2: 32 -> 64 通道
        self.conv3 = nn.Conv2d(32, 64, kernel_size=3,
                               padding=1, bias=False)
        self.bn3 = nn.BatchNorm2d(64)
        self.conv4 = nn.Conv2d(64, 64, kernel_size=3,
                               padding=1, bias=False)
        self.bn4 = nn.BatchNorm2d(64)
        self.pool2 = nn.MaxPool2d(2, 2)
        self.dropout2 = nn.Dropout2d(0.2)

        # Block 3: 64 -> 128 通道
        self.conv5 = nn.Conv2d(64, 128, kernel_size=3,
                               padding=1, bias=False)
        self.bn5 = nn.BatchNorm2d(128)
        self.conv6 = nn.Conv2d(128, 128, kernel_size=3,
                               padding=1, bias=False)
        self.bn6 = nn.BatchNorm2d(128)
        self.pool3 = nn.MaxPool2d(2, 2)
        self.dropout3 = nn.Dropout2d(0.3)

        # 分类�?         # 经过3次pooling (32 -> 16 -> 8 -> 4)
        # 特征维度: 128 * 4 * 4 = 2048
        self.classifier = nn.Sequential(
            nn.Linear(128 * 4 * 4, 256),
            nn.ReLU(inplace=True),
            nn.Dropout(0.4),
            nn.Linear(256, num_classes)
        )

        self._init_weights()

    def _init_weights(self):
        for m in self.modules():
            if isinstance(m, nn.Conv2d):
                nn.init.kaiming_normal_(
                    m.weight, mode='fan_out',
                    nonlinearity='relu'
                )
            elif isinstance(m, nn.BatchNorm2d):
                nn.init.constant_(m.weight, 1)
                nn.init.constant_(m.bias, 0)
            elif isinstance(m, nn.Linear):
                nn.init.kaiming_normal_(
                    m.weight, mode='fan_out',
                    nonlinearity='relu'
                )

    def forward(self, x):
        # Block 1
        x = F.relu(self.bn1(self.conv1(x)))
        x = F.relu(self.bn2(self.conv2(x)))
        x = self.pool1(x)
        x = self.dropout1(x)

        # Block 2
        x = F.relu(self.bn3(self.conv3(x)))
        x = F.relu(self.bn4(self.conv4(x)))
        x = self.pool2(x)
        x = self.dropout2(x)

        # Block 3
        x = F.relu(self.bn5(self.conv5(x)))
        x = F.relu(self.bn6(self.conv6(x)))
        x = self.pool3(x)
        x = self.dropout3(x)

        # 分类�?        x = x.view(x.size(0), -1)
        x = self.classifier(x)
        return x

# ========== 实例化模型 ==========
device = torch.device('cuda' if torch.cuda.is_available() else 'cpu')
model = CIFAR10CNN().to(device)
criterion = nn.CrossEntropyLoss()
optimizer = optim.SGD(model.parameters(), lr=0.1,
                      momentum=0.9, weight_decay=5e-4)
scheduler = optim.lr_scheduler.CosineAnnealingLR(
    optimizer, T_max=200
)

print(f"Model parameters: "
      f"{sum(p.numel() for p in model.parameters()):,}")

# ========== 训练函数 ==========
def train_epoch(model, loader, criterion, optimizer, device):
    model.train()
    total_loss, correct = 0, 0
    for X, y in tqdm(loader, desc='Train', leave=False):
        X, y = X.to(device), y.to(device)
        optimizer.zero_grad()
        output = model(X)
        loss = criterion(output, y)
        loss.backward()
        optimizer.step()
        total_loss += loss.item() * X.size(0)
        correct += (output.argmax(1) == y).sum().item()
    return (total_loss / len(loader.dataset),
            correct / len(loader.dataset))

@torch.no_grad()
def evaluate(model, loader, criterion, device):
    model.eval()
    total_loss, correct = 0, 0
    all_preds, all_labels = [], []
    for X, y in loader:
        X, y = X.to(device), y.to(device)
        output = model(X)
        loss = criterion(output, y)
        total_loss += loss.item() * X.size(0)
        pred = output.argmax(1)
        correct += (pred == y).sum().item()
        all_preds.extend(pred.cpu().tolist())
        all_labels.extend(y.cpu().tolist())
    return (total_loss / len(loader.dataset),
            correct / len(loader.dataset),
            all_preds, all_labels)

# ========== 训练循环 ==========
n_epochs = 80
best_acc = 0
history = {'train_loss': [], 'train_acc': [],
           'test_loss': [], 'test_acc': []}

for epoch in range(1, n_epochs + 1):
    train_loss, train_acc = train_epoch(
        model, train_loader, criterion, optimizer, device
    )
    test_loss, test_acc, _, _ = evaluate(
        model, test_loader, criterion, device
    )
    scheduler.step()

    history['train_loss'].append(train_loss)
    history['train_acc'].append(train_acc)
    history['test_loss'].append(test_loss)
    history['test_acc'].append(test_acc)

    if test_acc > best_acc:
        best_acc = test_acc
        torch.save(model.state_dict(), 'best_cifar10_cnn.pt')

    print(f"Epoch {epoch:3d}/{n_epochs} | "
          f"Train Loss: {train_loss:.4f} | "
          f"Train Acc: {train_acc:.4f} | "
          f"Test Loss: {test_loss:.4f} | "
          f"Test Acc: {test_acc:.4f}")

print(f"\nBest Test Accuracy: {best_acc:.4f}")

# ========== 加载最佳模型并测试 ==========
model.load_state_dict(torch.load('best_cifar10_cnn.pt'))
test_loss, test_acc, test_preds, test_labels = evaluate(
    model, test_loader, criterion, device
)
print(f"Final Test Accuracy: {test_acc:.4f}")
\end{lstlisting}
\end{mycode}

\subsection{感受野的堆叠增长}

一个重要的设计原则：两个 $3\times3$ 卷积堆叠的感受野等价于一个 $5\times5$ 卷积，但参数量更少（$2 \times 3^2 < 5^2$），且中间有非线性激活。这正是VGG网络的设计哲学。

\begin{figure}[H]\centering
\begin{tikzpicture}[scale=1.0]
% Input
\foreach \i in {0,...,6} {
    \foreach \j in {0,...,6} {
        \fill[blue!5, draw=gray, thin] (\i*0.5, 3.0-\j*0.5) rectangle ++(0.5,0.5);
    }
}
\node[font=\footnotesize] at (1.5,-0.5) {输入（$7\times7$）};

% After first 3x3 conv
\draw[->, thick] (3.7,1.5) -- (4.5,1.5);
\foreach \i in {0,...,4} {
    \foreach \j in {0,...,4} {
        \fill[green!15, draw=gray, thin] (\i*0.5+5.0, 2.5-\j*0.5) rectangle ++(0.5,0.5);
    }
}
\draw[red!70, very thick] (5.0, 2.5) rectangle (6.5, 1.0);
\node[font=\footnotesize] at (5.5,-0.5) {$3\times3$ 卷积后（$5\times5$）};

% After second 3x3 conv
\draw[->, thick] (7.7,1.5) -- (8.5,1.5);
\foreach \i in {0,...,2} {
    \foreach \j in {0,...,2} {
        \fill[orange!15, draw=gray, thin] (\i*0.5+9.0, 2.5-\j*0.5) rectangle ++(0.5,0.5);
    }
}
\draw[blue!70, very thick] (9.0, 2.5) rectangle (10.5, 1.0);
\node[font=\footnotesize] at (9.5,-0.5) {第2个$3\times3$卷积后};

\node[font=\small\bfseries, red!70] at (3.2, 3.8) {感受野 = 3};
\node[font=\small\bfseries, red!70] at (7.0, 3.8) {感受野 = 5};
\end{tikzpicture}
\caption{两个$3\times3$卷积堆叠的感受野增长。第一个卷积的感受野为3，第二个卷积作用在第一个的特征图上，感受野扩展为5。参数量为$2 \times 3^2 C^2$，而单个$5\times5$卷积需要$5^2 C^2$个参数}
\label{fig:receptive-field}
\end{figure}

%%%%%%%%%%%%%%%%%%%%%%%%%%%%%%%%%%%%%%%%%%%%%%%%%%%%%%%%
\section{转置卷积——从低分辨率到高分辨率}

\subsection{为什么需要上采样}

在图像分割和图像生成任务中，我们需要将低分辨率的特征图上采样回原始分辨率。\textbf{转置卷积}（Transpose Convolution，也称反卷积）是实现上采样的主要方法之一。

\begin{figure}[H]\centering
\begin{tikzpicture}[scale=1.0]
% Input 2x2
\node[font=\small\bfseries] at (-1.5, 2.0) {$\mathbf{X}$};
\foreach \i in {0,1} {
    \foreach \j in {0,1} {
        \fill[blue!15, draw=gray] (\i*0.8-1.0, 1.0-\j*0.8) rectangle ++(0.8,0.8);
        \node[font=\footnotesize] at (\i*0.8-0.6, 1.4-\j*0.8) {$x_{\i\j}$};
    }
}

% Arrow
\draw[->, very thick] (0.0,0.0) -- (1.8,0.0) node[midway, above, font=\footnotesize] {转置卷积};

% Output 4x4
\node[font=\small\bfseries] at (5.5, 2.0) {$\mathbf{Y}$};
\foreach \i in {0,...,3} {
    \foreach \j in {0,...,3} {
        \fill[green!15, draw=gray] (\i*0.8+2.8, 1.0-\j*0.8) rectangle ++(0.8,0.8);
        \node[font=\tiny] at (\i*0.8+3.2, 1.4-\j*0.8) {$y_{\i\j}$};
    }
}

\node[font=\footnotesize] at (3.0,-1.5) {每个输入元素``广播''到输出中对应的$k\times k$区域};
\end{tikzpicture}
\caption{转置卷积的概念。$2\times2$输入经$3\times3$转置卷积核上采样到$4\times4$（stride=2, padding=1）}
\label{fig:transposed-conv}
\end{figure}

\subsection{转置卷积的数学原理}

转置卷积是标准卷积的``反向''操作：标准卷积将输入展平为列向量 $\mathbf{x}$，卷积核展为行矩阵 $\mathbf{C}$，输出 $\mathbf{y} = \mathbf{C}\mathbf{x}$。转置卷积使用 $\mathbf{C}^{\top}$ 作为变换矩阵，将低分辨率特征图上采样：

\[
\mathbf{y}_{\text{trans}} = \mathbf{C}^{\top}\mathbf{x}_{\text{low}}
\]

\begin{mycode}{PyTorch 转置卷积示例}
\begin{lstlisting}[language=Python]
import torch
import torch.nn as nn

# 标准卷积：下采�?conv = nn.Conv2d(3, 64, kernel_size=3, stride=2, padding=1)
x = torch.randn(1, 3, 64, 64)
out = conv(x)
print(f"标准卷积: {x.shape} -> {out.shape}")  # (1,3,64,64) -> (1,64,32,32)

# 转置卷积：上采�?tconv = nn.ConvTranspose2d(64, 3, kernel_size=4,
                               stride=2, padding=1)
x_up = tconv(out)
print(f"转置卷积: {out.shape} -> {x_up.shape}")  # (1,64,32,32) -> (1,3,64,64)

# U-Net 风格的上采样模块
class UpBlock(nn.Module):
    def __init__(self, in_channels, out_channels):
        super().__init__()
        self.up = nn.ConvTranspose2d(
            in_channels, out_channels,
            kernel_size=2, stride=2
        )
        self.conv = nn.Sequential(
            nn.Conv2d(out_channels * 2, out_channels, 3,
                      padding=1),
            nn.BatchNorm2d(out_channels),
            nn.ReLU(inplace=True)
        )

    def forward(self, x, skip):
        x = self.up(x)
        x = torch.cat([x, skip], dim=1)  # 跳跃连接
        return self.conv(x)
\end{lstlisting}
\end{mycode}

%%%%%%%%%%%%%%%%%%%%%%%%%%%%%%%%%%%%%%%%%%%%%%%%%%%%%%%%
\section{特征可视化——理解CNN看到了什么}

\subsection{卷积核可视化}

第一层卷积核通常学习类似Gabor滤波器的边缘和颜色检测器：

\begin{mycode}{卷积核可视化}
\begin{lstlisting}[language=Python]
import torch
import torch.nn as nn
import torchvision
import matplotlib.pyplot as plt
import numpy as np

def visualize_kernels(model, layer_name='conv1'):
    """可视化指定卷积层的所有卷积核"""
    layer = dict(model.named_modules())[layer_name]
    kernels = layer.weight.data.cpu().numpy()

    # kernels shape: (C_out, C_in, kH, kW)
    C_out, C_in, kH, kW = kernels.shape

    # 归一化到 [0,1]
    kernels = (kernels - kernels.min()) / \
              (kernels.max() - kernels.min() + 1e-8)

    # 显示前 64 个卷积核
    n_display = min(C_out, 64)
    n_cols = 8
    n_rows = (n_display + n_cols - 1) // n_cols

    fig, axes = plt.subplots(n_rows, n_cols,
                             figsize=(12, 2 * n_rows))
    axes = axes.flatten()

    for i in range(n_display):
        # 显示第一个输入通道的卷积核
        if C_in == 3:
            axes[i].imshow(kernels[i].transpose(1, 2, 0))
        else:
            axes[i].imshow(kernels[i, 0], cmap='gray')
        axes[i].axis('off')

    for i in range(n_display, len(axes)):
        axes[i].axis('off')

    plt.suptitle(f'{layer_name} 卷积核可视化 '
                 f'(共{C_out}个，显示{n_display}个)',
                 fontsize=14)
    plt.tight_layout()
    plt.show()

# 使用示例
# model = trained_cnn_model
# visualize_kernels(model, 'conv1')
\end{lstlisting}
\end{mycode}

\subsection{特征图可视化}

观察中间层的激活模式，了解网络在不同深度提取了什么样的特征：

\begin{mycode}{特征图可视化}
\begin{lstlisting}[language=Python]
import torch
import torch.nn as nn
import matplotlib.pyplot as plt

class FeatureExtractor:
    """提取并可视化CNN中间层的特征图"""
    def __init__(self, model, layer_names):
        self.model = model
        self.features = {}
        self.hooks = []

        for name in layer_names:
            layer = dict(model.named_modules())[name]
            hook = layer.register_forward_hook(
                self._hook_fn(name)
            )
            self.hooks.append(hook)

    def _hook_fn(self, name):
        def hook(module, input, output):
            self.features[name] = output.detach()
        return hook

    def remove_hooks(self):
        for hook in self.hooks:
            hook.remove()

    def visualize(self, image, layer_name, max_channels=16):
        self.model.eval()
        with torch.no_grad():
            _ = self.model(image.unsqueeze(0))

        feat = self.features[layer_name][0]  # (C, H, W)
        C = min(feat.shape[0], max_channels)

        n_cols = 4
        n_rows = (C + n_cols - 1) // n_cols
        fig, axes = plt.subplots(n_rows, n_cols,
                                 figsize=(12, 3 * n_rows))
        axes = axes.flatten()

        for i in range(C):
            axes[i].imshow(feat[i].cpu().numpy(),
                          cmap='viridis')
            axes[i].axis('off')
            axes[i].set_title(f'Ch {i}', fontsize=8)

        for i in range(C, len(axes)):
            axes[i].axis('off')

        plt.suptitle(f'{layer_name} 特征�?'
                     f'(shape: {feat.shape})', fontsize=14)
        plt.tight_layout()
        plt.show()

# 使用示例
# extractor = FeatureExtractor(model, ['conv1', 'conv3', 'conv5'])
# extractor.visualize(sample_image, 'conv3')
# extractor.remove_hooks()
\end{lstlisting}
\end{mycode}

\subsection{Grad-CAM——类激活映射}

Grad-CAM利用目标类别的梯度在最后一个卷积层的特征图上的加权平均，生成类别的空间定位热力图：

\begin{mycode}{Grad-CAM 实现}
\begin{lstlisting}[language=Python]
import torch
import torch.nn.functional as F
import cv2
import numpy as np

class GradCAM:
    def __init__(self, model, target_layer):
        self.model = model
        self.target_layer = target_layer
        self.gradients = None
        self.activations = None

        # 注册 forward hook 获取激活�?        target_layer.register_forward_hook(
            self._save_activation
        )
        # 注册 backward hook 获取梯�?        target_layer.register_backward_hook(
            self._save_gradient
        )

    def _save_activation(self, module, inp, out):
        self.activations = out.detach()

    def _save_gradient(self, module, grad_in, grad_out):
        self.gradients = grad_out[0].detach()

    def generate(self, input_image, target_class=None):
        self.model.eval()
        output = self.model(input_image)

        if target_class is None:
            target_class = output.argmax(1).item()

        # 反向传播获取梯�?        self.model.zero_grad()
        one_hot = torch.zeros_like(output)
        one_hot[0, target_class] = 1
        output.backward(gradient=one_hot, retain_graph=True)

        # 全局平均池化梯度得到权�?        weights = self.gradients.mean(dim=(2, 3),
                                       keepdim=True)  # (1, C, 1, 1)

        # 加权组合特征�?        cam = (weights * self.activations).sum(dim=1,
                                                keepdim=True)
        cam = F.relu(cam)  # 只关注对类别有正向影响的区域

        # 归一化
        cam = cam - cam.min()
        cam = cam / (cam.max() + 1e-8)

        return cam.squeeze().cpu().numpy()

    def overlay(self, input_image, cam, alpha=0.5):
        """将热力图叠加到原图上"""
        cam = np.uint8(255 * cam)
        cam = cv2.applyColorMap(cam, cv2.COLORMAP_JET)

        img = input_image.squeeze().permute(1, 2, 0)
        img = (img - img.min()) / (img.max() - img.min() + 1e-8)
        img = np.uint8(255 * img.cpu().numpy())

        cam_resized = cv2.resize(cam, (img.shape[1],
                                       img.shape[0]))
        overlayed = cv2.addWeighted(img, 1 - alpha,
                                    cam_resized, alpha, 0)
        return overlayed

# 使用示例
# grad_cam = GradCAM(model, model.conv6)
# cam = grad_cam.generate(sample_image)
# overlay = grad_cam.overlay(sample_image, cam)
# plt.imshow(overlay)
\end{lstlisting}
\end{mycode}

%%%%%%%%%%%%%%%%%%%%%%%%%%%%%%%%%%%%%%%%%%%%%%%%%%%%%%%%
\section{CNN 设计中的核心原则}

\subsection{感受野的指数增长}

每层池化将空间分辨率减半，但卷积核的感受野在输入图上指数级增长。这是CNN能够捕捉大范围语义依赖的关键。

\[
\text{感受野}_l = \text{感受野}_{l-1} + (k_l - 1) \times \prod_{i=1}^{l-1} s_i
\]

\subsection{通道数递增，空间分辨率递减}

现代CNN的设计遵循一个经典模式：
\begin{itemize}
    \item 随深度增加，空间分辨率下降（通过池化或步长为2的卷积）
    \item 随深度增加，通道数增加（从3到64、128、256、512）
\end{itemize}

这种设计将空间信息``交换''为语义信息——浅层的大尺寸特征图编码精确的空间位置（适合定位），深层的小尺寸特征图编码丰富的语义类别信息（适合分类）。

\begin{figure}[H]\centering
\begin{tikzpicture}[
    myblock/.style={rectangle, draw=blue!60, fill=blue!5, minimum width=3.0cm, minimum height=1.0cm, font=\footnotesize, rounded corners=2pt, align=center},
    myarrow/.style={->, >=stealth, thick, gray!70},
]
\node[myblock, fill=blue!15] (l1) at (0,4) {$3 \times 32 \times 32$ \\ 输入图像};
\node[myblock, fill=green!15] (l2) at (0,2) {$32 \times 32 \times 32$ \\ 低级特征（边缘）};
\node[myblock, fill=green!15] (l3) at (0,0) {$32 \times 16 \times 16$ \\ Pool1};
\node[myblock, fill=orange!15] (l4) at (0,-2) {$64 \times 16 \times 16$ \\ 中级特征（纹理）};
\node[myblock, fill=orange!15] (l5) at (0,-4) {$64 \times 8 \times 8$ \\ Pool2};
\node[myblock, fill=red!15] (l6) at (0,-6) {$128 \times 8 \times 8$ \\ 高级特征（部件）};
\node[myblock, fill=red!15] (l7) at (0,-8) {$128 \times 4 \times 4$ \\ Pool3};

\draw[myarrow] (l1) -- (l2);
\draw[myarrow] (l2) -- (l3);
\draw[myarrow] (l3) -- (l4);
\draw[myarrow] (l4) -- (l5);
\draw[myarrow] (l5) -- (l6);
\draw[myarrow] (l6) -- (l7);

\node[font=\small, anchor=west] at (4.5,4) {通道：3};
\node[font=\small, anchor=west] at (4.5,2) {通道：32（细粒度空间）};
\node[font=\small, anchor=west] at (4.5,0) {通道：32（空间减半）};
\node[font=\small, anchor=west] at (4.5,-2) {通道：64（纹理组合）};
\node[font=\small, anchor=west] at (4.5,-4) {通道：64（空间再减半）};
\node[font=\small, anchor=west] at (4.5,-6) {通道：128（语义抽象）};
\node[font=\small, anchor=west] at (4.5,-8) {通道：128（最小空间）};

\draw[decorate, decoration={brace, mirror, raise=3pt}, very thick] (-2.5,4.5) -- (-2.5,-8.5) node[midway, left, font=\small] {深度};
\end{tikzpicture}
\caption{CNN的特征金字塔。空间分辨率逐层降低，通道宽度逐层增加，从空间精确的低级特征过渡到语义丰富的高级特征}
\label{fig:feature-pyramid}
\end{figure}

\subsection{$1\times1$ 卷积——通道间的信息交换}

$1\times1$卷积看起来像一个点（没有空间范围），但它在通道维度上执行完整的线性变换：

\[
\mathbf{Y}[o, i, j] = \sum_{c} K[o, c] \cdot \mathbf{X}[c, i, j] + b[o]
\]

$1\times1$卷积的用途：
\begin{itemize}
    \item \textbf{改变通道数}：不改变空间尺寸的情况下调整通道维度（升降维）
    \item \textbf{跨通道信息融合}：学习通道之间的线性组合
    \item \textbf{减少计算量}：在$3\times3$或$5\times5$卷积前先用$1\times1$降维（bottleneck结构）
    \item \textbf{增加非线性}：配合ReLU，在不增加过多参数的情况下提升网络表达能力
\end{itemize}

\subsection{空洞卷积——扩大感受野而不增加参数}

\textbf{空洞卷积}（Dilated / Atrous Convolution）在标准卷积核的元素之间插入``空洞''（零值间隙），使得在不增加参数量和计算量的情况下成倍扩大感受野：

\[
\mathbf{Y}[i, j] = \sum_{m=0}^{k-1}\sum_{n=0}^{k-1} \mathbf{X}[i + m\cdot d,\; j + n\cdot d] \cdot \mathbf{K}[m, n]
\]

其中 $d$ 是空洞率（dilation rate），$d=1$ 退化为标准卷积，$d=2$ 时 $3\times3$ 卷积的感受野等效为 $5\times5$。

\begin{figure}[H]\centering
\begin{tikzpicture}[scale=0.85]
% Standard convolution (d=1)
\node[font=\small\bfseries] at (-4,3.5) {$d=1$（标准卷积）};
\foreach \i in {0,...,6} {
    \foreach \j in {0,...,6} {
        \fill[blue!5, draw=gray, thin] (\i*0.6 -5.5, 3.2-\j*0.6) rectangle ++(0.6,0.6);
    }
}
\foreach \i in {2,3,4} {
    \foreach \j in {2,3,4} {
        \fill[red!40] (\i*0.6 -5.5, 3.2-\j*0.6) rectangle ++(0.6,0.6);
        \node[font=\tiny, white] at (\i*0.6 -5.2, 3.5-\j*0.6) {$\times$};
    }
}
\draw[red!70, very thick] (-4.3, 2.0) rectangle (-2.5, 0.2) node[midway, font=\tiny, white] {};

% Arrow
\draw[->, thick] (-1.0, 1.5) -- (1.0, 1.5);
\node[font=\footnotesize] at (0,1.0) {感受野 = 3};

% Dilated convolution (d=2)
\node[font=\small\bfseries] at (5,3.5) {$d=2$（空洞卷积）};
\foreach \i in {0,...,6} {
    \foreach \j in {0,...,6} {
        \fill[blue!5, draw=gray, thin] (\i*0.6 +3.5, 3.2-\j*0.6) rectangle ++(0.6,0.6);
    }
}
\foreach \i in {1,3,5} {
    \foreach \j in {1,3,5} {
        \fill[red!40] (\i*0.6 +3.5, 3.2-\j*0.6) rectangle ++(0.6,0.6);
        \node[font=\tiny, white] at (\i*0.6 +3.8, 3.5-\j*0.6) {$\times$};
    }
}
\draw[red!70, very thick] (4.7, 2.6) rectangle (7.1,-0.4) node[midway, font=\tiny, white] {};

\draw[->, thick] (8.0, 1.5) -- (10.0, 1.5);
\node[font=\footnotesize] at (9.0,1.0) {感受野 = 5};
\end{tikzpicture}
\caption{空洞卷积对比。$d=1$时$3\times3$卷积的感受野为3；$d=2$时同尺寸卷积核的感受野扩展为5，但不增加参数数量和计算量。空洞卷积广泛应用于语义分割任务（如DeepLab系列）}
\label{fig:dilated-conv}
\end{figure}

空洞卷积在语义分割中尤为重要——它允许网络在不损失分辨率（不池化）的情况下获得大的感受野，同时利用多尺度的上下文信息。

%%%%%%%%%%%%%%%%%%%%%%%%%%%%%%%%%%%%%%%%%%%%%%%%%%%%%%%%
\section{实践建议与常见陷阱}

\begin{itemize}
    \item \textbf{$3\times3$是标准核大小}——$5\times5$和$7\times7$通常可用多个$3\times3$堆叠替代，更加高效
    \item \textbf{BatchNorm几乎总是有益}——在CNN中加入BN可显著加速收敛并允许更大学习率
    \item \textbf{用卷积替代全连接}——全连接层可用$1\times1$卷积或全局平均池化替代，减少参数量
    \item \textbf{数据增强是免费的午餐}——对于小数据集，数据增强（翻转、裁剪、颜色抖动）的提升往往大于改进模型架构
    \item \textbf{检查感受野覆盖范围}——确保最深层的感受野覆盖输入图像的大部分区域
    \item \textbf{注意内存占用}——中间特征图的存储是显存的主要消费者，可使用梯度检查点技术trade计算换内存
\end{itemize}

\subsection{计算卷积参数的公式}

给定输入尺寸 $C_{\text{in}} \times H \times W$，卷积核 $C_{\text{out}} \times C_{\text{in}} \times k \times k$：

\[
\text{参数量} = C_{\text{out}} \times (C_{\text{in}} \times k^2 + 1) \quad (\text{如果 bias=True})
\]

\[
\text{FLOPs} \approx 2 \times C_{\text{out}} \times C_{\text{in}} \times k^2 \times H_{\text{out}} \times W_{\text{out}}
\]

%%%%%%%%%%%%%%%%%%%%%%%%%%%%%%%%%%%%%%%%%%%%%%%%%%%%%%%%
\section{本章小结}

本章是全书代码最密集的一章，我们深入探讨了卷积神经网络从理论到实践的每一个层面：

\begin{itemize}
    \item 卷积利用\textbf{局部连接}和\textbf{参数共享}将视觉任务的参数需求从 $O(H^2W^2)$ 降至 $O(k^2)$
    \item 用NumPy亲手实现的conv2d和im2col揭示了卷积的底层计算模式
    \item 池化层通过降采样增大感受野并引入局部平移不变性
    \item 完整的PyTorch CNN训练管线——数据增强、模块化设计、训练循环和学习率调度
    \item 转置卷积实现了从低分辨率特征到高分辨率输出的上采样
    \item 特征可视化（卷积核、特征图、Grad-CAM）是理解CNN内部表示的重要工具
    \item CNN设计的核心模式：空间分辨率递减、通道数递增，将空间信息转化为语义信息
\end{itemize}
